\documentclass[11pt]{article}
\usepackage[T1]{fontenc}
\usepackage{lmodern,amsmath,amssymb,amsthm,mathtools}
\usepackage[margin=1in]{geometry}
\usepackage{microtype}
\usepackage[colorlinks=true,allcolors=blue]{hyperref}
\newcommand{\T}{\mathbb T}
\newcommand{\R}{\mathbb R}
\newcommand{\eps}{\varepsilon}
\newcommand{\norm}[1]{\left\|#1\right\|}
\newcommand{\dist}{\operatorname{dist}}
\newtheorem{theorem}{Theorem}[section]
\newtheorem{lemma}[theorem]{Lemma}
\newtheorem{proposition}[theorem]{Proposition}
\newtheorem{corollary}[theorem]{Corollary}
\theoremstyle{remark}\newtheorem{remark}[theorem]{Remark}
\numberwithin{equation}{section}
\setlength{\emergencystretch}{2em}
\hypersetup{pdfauthor={Myeongjin Kang},pdftitle={Second-order contact of a ball and a slab in the cubic three-torus}}
\title{Second-order contact of a ball and a slab\texorpdfstring{\\}{ }in the cubic three-torus}
\author{Myeongjin Kang\thanks{Copyright 2026 Myeongjin Kang. Original paper and prose licensed under \href{https://creativecommons.org/licenses/by/4.0/}{CC BY 4.0}.}}
\date{October 14, 2026\quad---\quad Draft for independent manuscript review}
\begin{document}
\maketitle
\begin{abstract}
Relative to the pinned scalar isoperimetric classification of the cubic
three-torus, we establish a second-order contact law for two nested sets
whose volumes approach those of a tangent ball and slab. The admissible
sets have arbitrary finite perimeter. If the ball radius approaches the
ball--tube transition from below, with no restriction on its rate relative
to the compression, the ball-localized minimum is
$2\pi\eps^2/(|\log\eps|+\log|\log\eps|)
+O(\eps^2/|\log\eps|^2)$, uniformly in both parameters.
The lower bound uses the full positive spherical Jacobi inverse and the
two face Laplace inverses on actual deformed circles. A summed trace cancels
relative axial displacement, and improved convergence of quasiminimizers
provides a quantitative nonlinear error by interpolation. A tangent
logarithmic profile with exact volume corrections gives the matching upper
bound. The exact competing tube minimum then yields a second-scale
crossover law. The logarithmic inversion is classical; the assertion here
concerns the unrestricted nonlinear minimum and its uniform error.
\end{abstract}

\section{Statement and dependencies}\label{sec:statement}
Let $P$ denote De Giorgi perimeter on the unit cubic flat torus
$\T^3=\R^3/\mathbb Z^3$. All inclusions are understood almost everywhere.
For $R<1/3$ and small $\eps>0$, set
\begin{align}
 D_R(E,F)&=P(E)+P(F)-4\pi R^2-2,\\
 L_R(\eps)&=\inf\{D_R(E,F):E\subset F,\quad
 |E|=4\pi R^3/3,\quad |F|=2(R-\eps)\}.\label{eq:global}
\end{align}
The infimum is over all finite-perimeter pairs. There is no symmetry,
smoothness, topology, or connectedness assumption. Define
\[
 R_*=\frac13,\quad \rho_*=\frac29,\quad a_*=\frac{4\pi}{81},
 \quad \ell=|\log\eps|,\quad S=\frac{\eps^2}{\ell},
 \quad T_2=\frac{\eps^2\log\ell}{\ell^2}.
\]
Let $\mathcal B_*$ be the translation orbit of radius-$R_*$ balls and put
$\eta_0=a_*/10$. Let $M_B(R,\eps)$ be the infimum in
\eqref{eq:global} with the additional closed constraint
$\dist_{L^1}(E,\mathcal B_*)\leq\eta_0$. There is no additional
constraint on $F$.

\begin{theorem}\label{thm:main}
Uniformly as $R<1/3$ tends to $1/3$ and $\eps\downarrow0$,
\begin{equation}
 M_B(R,\eps)=\frac{2\pi\eps^2}{\ell+\log\ell}
                   +O\left(\frac{\eps^2}{\ell^2}\right).
 \label{eq:main}
\end{equation}
For all sufficiently close parameters,
\begin{equation}
 L_R(\eps)=\min\{M_B(R,\eps),g(R)\},\qquad
 g(R)=2\pi\sqrt{4R^3/3}-4\pi R^2.\label{eq:two-well}
\end{equation}
Consequently, for every family $R=R_\eps<1/3$ tending to $1/3$ such that
\[
 \frac{1/3-R_\eps-3S}{T_2}\longrightarrow\alpha\in\R,
\]
one has
\begin{equation}
 \frac{L_{R_\eps}(\eps)-2\pi S}{T_2}
 \longrightarrow\min\{-2\pi,(2\pi/3)\alpha\}.
 \label{eq:crossover}
\end{equation}
\end{theorem}
Uniformity means that one error constant and one parameter neighborhood
work in \eqref{eq:main}; no relation between $1/3-R$ and $\eps$ is imposed.
The crossover does not assert uniqueness of a finite-parameter transition,
coexistence of minimizers, or circularity of an unknown contact set.

We use the following precise scalar premise from \cite[Theorem 1.1,
label \texttt{intro:main}]{OpenAI}, at the immutable revision recorded
in the bibliography:
\begin{equation}
 I(V)=\min\{(36\pi)^{1/3}v^{2/3},2\sqrt{\pi v},2\},
 \qquad v=\min(V,1-V).\label{eq:profile}
\end{equation}
Its exhaustive equality classification consists of balls, tubes about
shortest coordinate geodesics, slabs, and their complements, with precisely
the adjacent types at transition volumes. We treat that theorem as a
premise, rather than re-proving it. The external regularity input is
\cite[Definition 2.5 and Theorem 4.2]{AFM}: uniformly quasiminimizing sets
with bounded perimeters that converge in $L^1$ to a smooth set have
whole-boundary normal graphs converging in $C^{1,\beta}$ for every
$\beta<1/2$. We verify its hypotheses below.

The leading ball--slab law and the localization/tube comparison are
antecedents of this work. Their needed arguments are reproduced here, so
no internal review verdict is a mathematical premise. Classical compressed
droplet calculations already contain the contact logarithm and its
log--log inversion; see \cite[Section III]{LGL} and
\cite[Sections 4.1--4.2 and Appendix B.2]{Lv}. A formal radial calculation
does not establish the lower bound over the class in \eqref{eq:global}.
The additional assertion in \eqref{eq:main} is its quantitative nonlinear
lower bound with uniform error. We make no exhaustive priority claim.

\section{Exact minimizing replacements}\label{sec:replacement}
All constants in this section are uniform for $R$ in a sufficiently small
closed interval ending at $R_*$. The outer volumes lie in
$(1/\pi,1-1/\pi)$, so \eqref{eq:profile}, including complementation,
gives $P(F)\geq2$. For $R\leq R_*$ the inner scalar bound is
$P(E)\geq4\pi R^2$. The two deficits are separately nonnegative.

\begin{lemma}\label{lem:upper-rough}
There are admissible ball-neighborhood pairs of deficit $O(S)$, uniformly
in the two parameters. There are also admissible tube--slab pairs of
deficit $g(R)$ with a uniform positive clearance between their boundaries.
\end{lemma}
\begin{proof}
For the second assertion put $\rho=\sqrt{4R^3/3}$ and take the tube
$\{y^2+z^2<\rho^2\}$ with coordinate axis $x$, inside the slab
$\{|z|<R-\eps\}$. The clearance tends to $R_*-\rho_*=1/9$.

For the first assertion fix a small polar chart radius $d$ and a fixed
$\gamma>0$. Put $a=\sqrt\eps\,\ell$ and
$H=\log(d/a)$. A radial cutoff $\chi$, equal to one for $r\leq a$
and to $\log(d/r)/H$ on the annulus, can be rounded at its endpoints so that
\[
 \int|D\chi|^2=2\pi/H+O(H^{-2}),\quad
 \int\chi^j=O_d(H^{-j})+O(a^2)\quad(j=1,2),
 \quad\norm{D\chi}_\infty=O((aH)^{-1}).
\]
Move each spherical polar patch inward normally by
$t\chi$, $t=(1+\gamma)\eps/2$, and each corresponding slab face
outward by $t\chi$. Fixed free equatorial and exterior-face patches
restore the two volumes exactly, with parameters $O_d(\eps/\ell)$.
Their volume derivatives are nonzero and their clearances are positive,
uniformly in $R$; the inverse function theorem applies.
For actual $|x|\leq a/2$, the spherical cutoff is still one after
normal-to-vertical conversion, and the total gain is
$(1+\gamma+o(1))\eps$. Outside this core the reference gap is at least
$ca^2\gg\eps$, until the fixed polar chart ends; farther away it is
uniformly positive. This pays all displacements and preserves inclusion.
Boundary containment implies region containment: the extrema of a
nonempty compact inner vertical slice occur on its boundary, and the
outer occupied slice is an interval. No connected-slice assumption is
needed. The exact-volume sphere area expansion and face expansions give
total cost $O(t^2/H)+O_d(\eps^2/\ell^2)=O(S)$. The unscaled
$C^1$ norms tend to zero. All charts, correction derivatives and constants
are uniform in the indicated compact radius interval.
\end{proof}

Let $\mathcal C_*$ be the union of the three coordinate-tube translation
orbits of radius $\rho_*$, and define $M_C$ using
$\dist_{L^1}(E,\mathcal C_*)\leq\eta_0$.
\begin{lemma}\label{lem:localize}
Both localized minima and the global minimum are attained. For all
sufficiently close parameters the localized constraints are inactive at
their minimizers, and $L_R=\min(M_B,M_C)$.
\end{lemma}
\begin{proof}
BV compactness, lower semicontinuity of perimeter, and closure of inclusion
and the two volume constraints give attainment. The orbit-distance
constraints are $L^1$ closed. A radius-$R_*$ ball projects on any
coordinate circle into an interval of length $2/3$. Along its own axis
a radius-$\rho_*$ tube has constant cross-sectional area $a_*$.
Thus $|C\setminus B|\geq a_*/3$, and equal volumes imply
$|C\mathbin\triangle B|\geq2a_*/3$. Their $\eta_0$ neighborhoods
are disjoint.

Any sequence of admissible pairs with parameters tending to $(R_*,0)$
and deficit tending to zero has scalar equality limits, since both
deficits are nonnegative. By \eqref{eq:profile}, the inner limit is a
ball or coordinate tube and the outer limit an occupied slab of width
$2/3$. Inclusion persists. A tube normal to the slab is impossible;
a parallel one may have normal center offset at most $1/9$. A ball's
projection has exactly the slab width, so their normal centers agree.
The upper bounds in Lemma~\ref{lem:upper-rough} tend to zero in either
localized class, even if $g(R)/S$ diverges. Separation forces every
localized limit into its own orbit, making its distance tend to zero.
The constraint is therefore interior by a fixed margin. Global exact
minimizers have the ball upper bound $O(S)$; their limits lie in one of
the two orbits. A counterexample sequence to the asserted identity would
contradict this compactness conclusion.
\end{proof}

\begin{lemma}\label{lem:graphs}
Along every converging normalized sequence of localized exact minimizers,
each scalar boundary is a whole-boundary $C^{1,\beta}$ normal graph over
its smooth limiting boundary, with norm tending to zero, for $\beta<1/2$.
The bounds remain uniform after changing to the parameter-dependent
reference radius and slab width and after reference centering.
\end{lemma}
\begin{proof}
At every limiting nested pair choose an inner boundary patch strictly
inside the slab and an outer face patch strictly outside the inner set.
Such patches also exist at tube--face tangency. Two smooth ambient vector
fields supported there have a diagonal invertible two-volume derivative.
For nearby measurable pairs the volume maps are integrals of their smooth
flow Jacobians, so $L^1$ closeness preserves a uniform inverse bound.
A common ambient diffeomorphism corrects total volume error $q$ with
parameter $O(q)$, perimeter cost $O(q)$ for bounded-perimeter sets,
and $L^1$ displacement $O(q)$. The last estimate follows by smooth
approximation and the BV transport bound along the flow. Inclusion is
preserved by the common map.

For an arbitrary inner scalar comparison $H$, form
$(H\cap F,H\cup F)$. Perimeter submodularity bounds its sum by
$P(H)+P(F)$, and its total paired symmetric difference is exactly
$|H\mathbin\triangle E|$. This identity follows pointwise from the
three possible nested states $(0,0),(0,1),(1,1)$. For an outer comparison
use $(E\cap H,E\cup H)$, with difference $|H\mathbin\triangle F|$.
For smaller-perimeter comparisons in a sufficiently small fixed ball,
the volume correction stays inside the inactive localization constraint.
Local paired minimality therefore gives the scalar inequality
$P(E)\leq P(H)+C|E\mathbin\triangle H|$, and likewise for $F$.
Larger-perimeter comparisons satisfy it automatically. For larger balls
use the uniform perimeter bound. Thus, for every ball $B_r$ and every
finite-perimeter comparison supported compactly in it,
\[
 P(E)\leq P(H)+\omega r^3,
\]
with one $\omega$ for the sequence, and similarly for $F$.
This is exactly the AFM hypothesis, without a scalar volume constraint.
Their theorem applies separately to the inner boundary and the disconnected
union of the two outer faces. The occupied slab side is selected by
$L^1$ convergence. Whole-boundary convergence excludes additional components
of these exact replacements.

Changing reference radius and width is a smooth small reparametrization.
Sphere centering uses the three degree-one moments: their derivative at
zero is the nonsingular translation Gram matrix. The near-identity graph
coordinate maps and their inverses have uniform $C^{1,\beta}$ bounds.
Face reference changes are constant shifts and translations. Hence these
operations preserve a uniform $C^{1,\beta}$ bound. If joint uniformity
failed, compactness would give a normalized equality limit and the same
AFM conclusion along a counterexample sequence. This is a contradiction.
\end{proof}

\begin{lemma}\label{lem:tube}
For all sufficiently close parameters, $M_C(R,\eps)=g(R)$ exactly.
\end{lemma}
\begin{proof}
Use a coordinate tube reference of radius $\rho=\sqrt{4R^3/3}$ and
write its inner graph as
$(s,(\rho+u)\cos\theta,(\rho+u)\sin\theta)$.
Transverse reference centering removes the moments against
$\cos\theta,\sin\theta$. Exact volume and area are respectively
\[
 \int(\rho u+u^2/2)\,ds\,d\theta=0,\qquad
 \int\sqrt{(\rho+u)^2(1+u_s^2)+u_\theta^2}\,ds\,d\theta.
\]
Eliminating the linear term gives
\[
 P(E)-2\pi\rho=\frac12\int
 [u_s^2+\rho^{-2}(u_\theta^2-u^2)]\rho\,ds\,d\theta
 +o(1)\norm u_{H^1}^2.
\]
The Jacobi eigenvalues are $4\pi^2 k^2+(m^2-1)/\rho^2$.
The constant mode is quadratic by the volume equation; the only zero
modes are the removed transverse translations. The remaining gap near
$\rho_*$ is bounded below by a positive constant, since
$\min(4\pi^2-81/4,243/4)>0$. Absorbing the quadratic mean and
small remainder yields $P(E)\geq2\pi\rho+c\norm u_{H^1}^2$.
Together with $P(F)\geq2$ this gives $M_C\geq g(R)$; the explicit
tube competitor gives the reverse inequality. Sequential contradiction
proves the assertion throughout a parameter neighborhood, without a
restriction on $(1/3-R)/S$.
\end{proof}

\section{Quantitative control of the ball graphs}\label{sec:quantitative}
Consider ball-localized exact minimizers. Normalize a limiting slab normal
to the $z$ axis. On the reference sphere $S_R$ write the radial height
$\psi$, centered to have zero degree-one moments. Write the two outward
face heights as $V^+,V^-$ over a reference slab of width $2(R-\eps)$.
Its exact volume makes the two means sum to zero; a reference axial
translation changes them oppositely, so both are made zero. Retain the
relative axial center $h$ of the sphere. These are reference changes of
the same nested pair, not independent motions of its sets.

Exact inner volume gives
\begin{equation}
 \int_{S_R}\left(\psi+\frac{\psi^2}{R}
                  +\frac{\psi^3}{3R^2}\right)=0.\label{eq:volume}
\end{equation}
The exact radial area density relative to $dS_R$ is
$(1+\psi/R)^2\sqrt{1+|D\psi|^2/(1+\psi/R)^2}$.
Using \eqref{eq:volume} gives
\begin{align}
 P(E)-4\pi R^2&=\tfrac12 Q_R(\psi)+\mathcal R_B,
 &Q_R(\psi)&=\int_{S_R}(|D\psi|^2-2\psi^2/R^2),\label{eq:area}\\
 |\mathcal R_B|&\leq C(\norm\psi_\infty+
 \norm{D\psi}_\infty^2)\norm\psi_{H^1}^2.\notag
\end{align}
For each flat face the error beyond half its Dirichlet integral is bounded
by $C\norm{DV}_\infty^2\norm{DV}_2^2$.
If $b=|S_R|^{-1}\int_{S_R}\psi$ and $f=\psi-b$, then the volume equation yields
$|b|\leq C\norm f_2^2$ after absorbing $b^2$. The spherical eigenvalues
$[l(l+1)-2]/R^2$ for $l\geq2$ are uniformly coercive in $H^1$;
the degree-one modes vanish. The mean and the relative small remainder
can be absorbed. Face Poincare inequalities apply to the zero means.
Lemma~\ref{lem:upper-rough} and nonnegativity of the separate deficits give
\begin{equation}
 \norm\psi_{H^1}+\norm{V^+}_{H^1}+\norm{V^-}_{H^1}\leq C\sqrt S,
 \qquad |b|\leq CS.\label{eq:H1}
\end{equation}

\begin{lemma}\label{lem:rate}
For these replacements,
\begin{equation}
 D_R(E,F)=\tfrac12 Q+O(S^{11/10}+S^2),\qquad
 Q=Q_R(f)+\int|DV^+|^2+\int|DV^-|^2.\label{eq:quant}
\end{equation}
The error is $o(\eps^2/\ell^2)$, uniformly in the two parameters.
\end{lemma}
\begin{proof}
Fix $\beta=1/4$ in Lemma~\ref{lem:graphs}. A gradient with bounded
$\beta$-Hölder seminorm and maximum size $M$ has magnitude at least
$M/2$ on a surface disk of radius $cM^{1/\beta}$. Thus
\[
 \norm{Du}_2^2\geq cM^{2+2/\beta},\qquad
 \norm{Du}_\infty\leq C\norm{Du}_2^{\beta/(1+\beta)}.
\]
Uniform charts suffice; the nonsmall-$M$ case is excluded by the small
$L^2$ norm. Connectedness bounds the function supremum by its mean and
gradient supremum. Apply this separately to the sphere and each face,
using \eqref{eq:H1}, to obtain a total $C^1$ norm $O(S^{1/10})$.
Equation~\eqref{eq:area}, the face errors, and
$Q_R(\psi)=Q_R(f)+O(b^2)$ prove \eqref{eq:quant}.
Finally $S^{11/10}/(\eps^2/\ell^2)=\eps^{1/5}\ell^{9/10}\to0$,
and $S^2/(\eps^2/\ell^2)=\eps^2\to0$.
\end{proof}

\section{The full inverse kernels and actual-curve traces}\label{sec:trace}
Let $G_T$ be the mean-zero inverse kernel of $-\Delta$ on the unit
two-torus. In a fixed diagonal chart,
\begin{equation}
 G_T(x,y)=-\frac1{2\pi}\log|x-y|+O(1).\label{eq:torus-kernel}
\end{equation}
One obtains this by cutting off the planar fundamental solution and
solving the remaining smooth mean-zero error by its Fourier series.
The correcting coefficients decay rapidly, so its kernel is bounded and
continuous. Away from the diagonal $G_T$ is bounded.

Let $K_R$ be the inverse kernel of $-\Delta_{S_R}-2/R^2$ restricted
to degrees $l\geq2$. Relative to physical sphere area, the addition
formula gives
\begin{equation}
 K_R(R\omega,R\eta)=\sum_{l\geq2}
 \frac{(2l+1)P_l(\omega\cdot\eta)}{4\pi[l(l+1)-2]}.
 \label{eq:sphere-kernel}
\end{equation}
The factor $R^{-1}$ in each normalized spherical harmonic cancels
the $R^{-2}$ in the eigenvalue. The mean-zero sphere Laplace inverse is
$-(4\pi)^{-1}\log(1-\omega\cdot\eta)$ plus a constant: direct radial
differentiation gives the off-pole equation and unit pole flux.
Its coefficient difference from \eqref{eq:sphere-kernel}, for $l\geq2$,
is, apart from $1/(4\pi)$,
\[
 \frac{2(2l+1)}{l(l+1)[l(l+1)-2]}
 =\frac1{l-1}-\frac1l-\frac1{l+1}+\frac1{l+2}.
\]
The positive series sums to $2/3$. Since $|P_l|\leq1$, this difference
is a uniformly bounded continuous kernel. Removing degree one is a
bounded finite-rank term. Consequently
\begin{equation}
 K_R(x,y)=-\frac1{2\pi}\log|x-y|+H_R(x,y),\qquad |H_R|\leq C,
 \label{eq:jacobi-log}
\end{equation}
uniformly on the compact radius interval; it is bounded between opposite
polar patches. Changing unit chord distance to physical distance adds
only a bounded constant involving $\log R$.

\begin{lemma}\label{lem:curve-trace}
For a bounded-density $C^1$ curve measure $\sigma$ with finite logarithmic
self-integral and a $C^1$ function $u$ in the corresponding mean-zero
or degree-$\geq2$ space,
\begin{equation}
 \left|\int u\,d\sigma\right|^2
 \leq Q_{\rm surf}(u)\iint K(x,y)\,d\sigma(x)d\sigma(y).
 \label{eq:trace-dual}
\end{equation}
The sources used below satisfy this hypothesis uniformly after their
curve scale is factored out.
\end{lemma}
\begin{proof}
For smooth source densities spectral Cauchy--Schwarz proves the inequality;
the projection is already part of the inverse kernel. Mollify the curve
measure in surface charts with positive kernels at scale $t<a$, where
$a$ is its curve scale. When original points are separated by more than
a fixed multiple of $t$, mollification changes the logarithm by a bounded
amount. At separation $O(t)$ its absolute value is bounded by
$C+\log(1/t)$, dominated by $C+|\log|x-y||$ after increasing $C$.
For the bilipschitz parametrizations below this has integrable majorant
$C+|\log(a|e^{i\theta}-e^{i\varphi}|)|$ at fixed $a$.
Dominated convergence gives convergence of the double integrals; the
bounded continuous kernel parts cause no difficulty. Pairings converge
first for smooth $u$, then for $C^1$ functions by simultaneous $C^0$
and $H^1$ approximation and the finite-dimensional projections.
This proves \eqref{eq:trace-dual}; it is a curve trace, not an $H^1$
point evaluation.
\end{proof}

At either pole let $F_\pm$ denote the actual inverse horizontal
parametrization into $S_R$ and put $p(x)=\psi(F_\pm(x))$.
The actual outward vertical displacement of the sphere satisfies the
exact identity
\begin{equation}
 W^\pm(x)=A_\pm(x)\psi(F_\pm(x)),\qquad
 A_\pm(x)=\frac{2R+p(x)}{
 \sqrt{(R+p(x))^2-|x|^2}+\sqrt{R^2-|x|^2}}.\label{eq:actual}
\end{equation}
At the bottom, outward means downward. In reference coordinates $y$ the
forward horizontal map is $y\mapsto(1+\psi(q_\pm(y))/R)y$, whose
derivative is uniformly close to the identity. Thus the actual curves
$\theta\mapsto F_\pm(ae^{i\theta})$ are bilipschitz at scale $a$:
their physical chord distances are between
$ca|e^{i\theta}-e^{i\varphi}|$ and
$Ca|e^{i\theta}-e^{i\varphi}|$ for every angular pair.
They remain in separated opposite polar patches. The identity
$A(p,0)=1$ holds for every small $p$; smoothness in $p$ and $|x|^2$
therefore gives
\begin{equation}
 |A_\pm(ae^{i\theta})-1|\leq Ca^2.\label{eq:weight}
\end{equation}
There is no additive error proportional to the graph's $C^1$ norm.

On the direct sum quadratic space in \eqref{eq:quant}, define
\[
 \Lambda_a(f,V)=\sum_\pm\frac1{2\pi}\int_0^{2\pi}
 [V^\pm(ae^{i\theta})-A_\pm(ae^{i\theta})
                      f(F_\pm(ae^{i\theta}))],d\theta.
\]
Each of the two face sources has mass one. Each sphere source has mass
$1+O(a^2)$. Equations~\eqref{eq:torus-kernel}--\eqref{eq:weight}
and the integrability of $|\log|e^{i\theta}-e^{i\varphi}||$ show that
each of the four self integrals is
$(2\pi)^{-1}\log(1/a)+O(1)+O(a^2\log(1/a))$.
The opposite-pole sphere cross term is bounded. The two faces belong
to different quadratic summands and have no cross term. Hence
\begin{equation}
 \Lambda_a(f,V)^2\leq\left(\frac2\pi\log\frac1a+C\right)Q.
 \label{eq:four-traces}
\end{equation}
The curves and weights depend on the minimizing graph, but are fixed
when applying the linear trace inequality; all constants are uniform.

\section{The quantitative lower bound}\label{sec:lower}
Actual inclusion in the vertical charts reads
\begin{equation}
 V^\pm(x)-W^\pm(x)\geq\eps\pm h-q_R(|x|),\qquad
 q_R(r)=R-\sqrt{R^2-r^2}.\label{eq:obstacle}
\end{equation}
For the lower face this follows from
$h-\sqrt{R^2-|x|^2}-W^-(x)\geq-(R-\eps)-V^-(x)$;
in particular its center term has sign $-h$.
Choose the deterministic test radius
\[
 a=\sqrt{\eps/\ell},\qquad L=\log(1/a)=\tfrac12(\ell+\log\ell).
\]
This is not an assertion about the shape or radius of a minimizer's
contact set. Average both inequalities on this circle and add.
The two center terms cancel exactly. Since $\psi=f+b$, $|b|\leq CS$,
and $q_R(a)=O(\eps/\ell)$ uniformly, one obtains
\begin{equation}
 \Lambda_a(f,V)\geq2\eps-C\eps/\ell-CS>0.\label{eq:forced}
\end{equation}
No estimate on $h/\eps$ is required. Combining
\eqref{eq:four-traces}, \eqref{eq:forced} and Lemma~\ref{lem:rate},
\begin{align}
 M_B(R,\eps)&\geq
 \frac{(2\eps-C\eps/\ell-CS)^2}{2[(2/\pi)L+C]}
       -C(S^{11/10}+S^2)\notag\\
 &\geq\frac{2\pi\eps^2}{\ell+\log\ell}
                  -C\frac{\eps^2}{\ell^2}.\label{eq:lower}
\end{align}
Every discarded term has the displayed scale or smaller. Uniformity
along arbitrary normalized sequences implies joint uniformity by
contradiction and compactness. The lower estimate concerns exact minima
over the full localized finite-perimeter class. Every original competitor
has at least its corresponding minimum; arbitrary competitors were never
assumed to be graphs or quasiminimizers.

\section{A matching recovery with exact volumes}\label{sec:recovery}
Fix a small $d>0$ below the chart radii and free-patch clearances. Let
$a=a(R,\eps)$ be the small solution of
\begin{equation}
 a^2[\log(d/a)+1/2]=R\eps,
 \qquad H=\log(d/a),\qquad A=\frac{a^2}{R}
                         =\frac{\eps}{H+1/2}.\label{eq:tangent}
\end{equation}
The left side is strictly increasing on $(0,d)$, with derivative
$2a\log(d/a)$, so the small solution exists uniquely for small $\eps$.
Define the radial profile
\[
 Z(r)=\begin{cases}
 \eps-r^2/(2R),&0\leq r\leq a,\\
 A\log(d/r),&a<r\leq d/2,
 \end{cases}
\]
and multiply the logarithmic tail by a fixed smooth nonnegative cutoff
on $d/2<r<d$, vanishing smoothly at $d$. Values and first derivatives
match at $a$ by \eqref{eq:tangent}. The profile is $C^{1,1}$, sufficient
for finite perimeter, with no vertical rim. Beyond $a$ the difference
between the logarithmic tail and the parabola has derivative
$r/R-a^2/(Rr)\geq0$ and begins at zero. In the cutoff region the
parabolic obstacle is negative. Since $q_R(r)\geq r^2/(2R)$,
\begin{equation}
 Z(r)\geq\eps-q_R(r).\label{eq:profile-obstacle}
\end{equation}

Use actual horizontal coordinates. Replace the upper sphere height
$\sqrt{R^2-r^2}$ by $\sqrt{R^2-r^2}-Z/2$, and reflect at the lower
pole. Move each corresponding outer face outward by $Z/2$.
The actual gap is $q_R(r)-\eps+Z(r)\geq0$. Outside the caps it is
positive. The inner slice argument from Lemma~\ref{lem:upper-rough}
proves full containment. The lost inner volume and gained outer volume
are $O(A)$. Independent corrections on fixed free equatorial and exterior
face patches restore the exact volumes with parameters $O(A)$ by the
inverse function theorem. Their positive gaps preserve inclusion.

Direct radial integration, before the fixed cutoff modification, gives
\begin{align}
 J:=\int|DZ|^2
 &=\frac{\pi a^4}{2R^2}+\frac{2\pi a^4H}{R^2}
   =\frac{2\pi\eps^2(H+1/4)}{(H+1/2)^2},\label{eq:J}\\
 \int Z&=O_d(A),\qquad \int Z^2=O_d(A^2),\qquad
 \int r^2|DZ|^2=O_d(A^2),\qquad
 \norm{DZ}_\infty=O(\sqrt{\eps/\ell}).\label{eq:weighted}
\end{align}
The cutoff changes $J$ by $O_d(A^2)$ and preserves the other estimates.
For example its annular weighted-gradient integral is bounded by
$2\pi A^2\int_a^d r\,dr$; the core contributes $O(a^6)=o(A^2)$.
The substitution $s=\log(d/r)$ proves the annular $L^1,L^2$ bounds;
their core contributions are smaller. The defining equation already
gives $H\sim\ell/2$, so $A=O(\eps/\ell)$ uniformly.

It is essential to keep integrated curvature errors rather than a fixed
relative chart distortion. On the upper reference cap
$s_0(x)=\sqrt{R^2-|x|^2}$,
\[
 \frac{Ds_0}{\sqrt{1+|Ds_0|^2}}=-\frac{x}{R}.
\]
For the compactly supported vertical change $v=-Z/2$, integration by
parts gives the exact first variation $(2/R)\int v$, which is $2/R$
times its signed volume change. The reflected lower cap has the same
outward convention. The exact equatorial volume correction contributes
$2/R$ times the restored volume plus $O(A^2)$; its support is disjoint
from the caps. Thus all inner linear terms cancel.

The Hessian of $\sqrt{1+|p|^2}$ at $p=Ds_0$ is $I+O(r^2)$.
Its quadratic correction is bounded by
$C\int r^2|Dv|^2=O(A^2)$. The cubic remainder is bounded by
$C\norm{Dv}_\infty\int|Dv|^2$, which is
$o(\eps^2/\ell^2)$ because its ratio to that scale is
$O(\sqrt{\eps\ell})$. The correcting bump has area minus
$(2/R)$ times volume equal to $O(A^2)$ by smooth Taylor expansion.
Flat faces have zero linear area variation and quartic-gradient errors;
their volume correction costs $O(A^2)$ as well.
Each of four patches has height magnitude $Z/2$ and half-Dirichlet area
coefficient. Their total deficit is consequently
\begin{equation}
 D_R(\text{recovery})\leq\frac J2+C\frac{\eps^2}{\ell^2}
 =\frac{\pi\eps^2(H+1/4)}{(H+1/2)^2}
       +C\frac{\eps^2}{\ell^2}.\label{eq:recovery}
\end{equation}
Taking logarithms in \eqref{eq:tangent} yields
\[
 2H=\ell+\log(H+1/2)+\log(d^2/R),\qquad
 H=\tfrac12(\ell+\log\ell)+O_d(1).
\]
For clarity, the first relation gives $H\sim\ell/2$; substituting this
in its logarithm proves the second. All constants are uniform in $R$.
Equation~\eqref{eq:recovery} matches \eqref{eq:lower}. The constructed
inner sets converge to the reference balls and belong to the ball
neighborhood eventually. This proves \eqref{eq:main}.

\section{Crossover and scope of verification}\label{sec:conclusion}
Lemmas~\ref{lem:localize} and \ref{lem:tube} prove
\eqref{eq:two-well}. Expanding the reciprocal in \eqref{eq:main} gives
\[
 M_B=2\pi S-2\pi T_2+o(T_2),
\]
since the reciprocal remainder is
$O((\log\ell)^2/\ell^3)$ and $\eps^2/\ell^2=o(T_2)$.
Direct differentiation gives $g'(1/3)=-2\pi/3$, with bounded second
derivative near $1/3$. If $\delta=1/3-R=3S+\alpha T_2+o(T_2)$,
then $\delta=O(S)$ and
$\delta^2/T_2=O(\eps^2/\log\ell)\to0$. Hence
$(g(R)-2\pi S)/T_2\to(2\pi/3)\alpha$.
Taking the minimum of the two convergent normalized costs proves
\eqref{eq:crossover}, completing Theorem~\ref{thm:main}.

The proof uses the scalar classification and the established AFM
regularity theorem at the exact interfaces stated above. It does not
classify the contact set, prove a finite-parameter unique crossing, or
give regularity to arbitrary near-minimizers. The tangent radius in the
recovery is a construction parameter, not an asserted radius of an
unknown minimizer. The second-scale crossover is a consequence of one
contact theorem, not a separate result.

The accompanying exact-arithmetic program checks coefficient identities
for the inverse-kernel difference, the tangent-profile energy, and the
elementary crossover expansion. These finite checks do not certify
regularity, measure convergence, uniform asymptotics, or the theorem.
No Lean formalization of this geometric argument is claimed. The earlier
leading contact argument and the moving-radius localization are credited
antecedents reproduced in the proof. The new quantitative passage uses
full inverse kernels and actual deformed curve measures, with an error
small enough to resolve the logarithmic correction.

\begin{thebibliography}{9}
\raggedright
\bibitem{OpenAI} OpenAI, \emph{The Isoperimetric Conjecture for the Cubic
Flat Three-Torus}, September 24, 2026. Theorem \texttt{intro:main},
\href{https://github.com/openai/math/tree/fd4aeeb2ee4fc729c18d98444fed42fd0529eeeb}{immutable repository revision}
\texttt{fd4aeeb2ee4fc729c18d98444fed42fd0529eeeb}.
\bibitem{AFM} E. Acerbi, N. Fusco and M. Morini,
\emph{Minimality via second variation for a nonlocal isoperimetric problem},
Communications in Mathematical Physics \textbf{322} (2013), 515--557.
\href{https://cvgmt.sns.it/media/doc/paper/540/Acerbi-Fusco-Morini.FINAL.pdf}{Primary manuscript},
Definition 2.5 and Theorem 4.2.
\bibitem{LGL} M.-D. Lacasse, G. S. Grest and D. Levine,
\emph{Deformation of Small Compressed Droplets}, Physical Review E
\textbf{54} (1996), 5436.
\href{https://arxiv.org/abs/mtrl-th/9603005v1}{arXiv:mtrl-th/9603005v1}, Section III.
\bibitem{Lv} C. Lv, \emph{Effects of geometric confinement on a droplet
between two parallel planes},
\href{https://arxiv.org/abs/1704.05830v2}{arXiv:1704.05830v2} (2017).
\end{thebibliography}
\end{document}
